\subsection{Results}
\begin{frame}
    \frametitle{Energy Supplied}
    \begin{columns}
        \column[t]{5cm}
            \begin{itemize}
                \item Initial deficit of over 5 GWe-yr resulting from \gls{LWR} 
                      refeuling outages
                \item Scenarios 5 (MMR + VOYGR) and 6 (Xe-100 + VOYGR) don't 
                      ever meet the demand
                \item Scenario 4 (Xe-100 + MMR) best meets the demand
                \item Scenario 3 (Xe-100) overprodces by 0.06 GWe-yr
                \item Flawed methodology that doesn't account for 
                      capacity factors
            \end{itemize}
        \column[t]{5cm}
        \vspace{-0.8cm}
            \begin{figure}
                \centering
                \includegraphics[scale=0.4]{nogrowth_energy_after_2025.pdf}
                \caption{Annual average energy produced (GWe-yr) in Scenarios 
                2-7 compared with energy demand}
                \label{fig:energy}

        \end{figure}
    \end{columns}
\end{frame}

\begin{frame}
    \frametitle{Number of Reactors}
    \begin{columns}
        \column[t]{4.5cm}
            \begin{itemize}
                \item Scenario 2 (MMR) deploys the most reactors
                \item Scenario 3 (Xe-100) deploys the fewest reactors
                \item Deploying VOYGRs with Xe-100 results in 
                      about the same number of reactors as deploying 
                      just Xe-100s
            \end{itemize}
        \column[t]{5.5cm}
        \vspace{-0.5cm}
            \begin{figure}
                \centering
                \includegraphics[scale=0.32]{nogrowth_reactors.pdf}
                %\caption{Annual avergae mass of enriched uranium required to fuel
                %advanced reactors in Scenarios 2-7.}
                %\label{fig:reactors}
        \end{figure}
        \vspace{-0.5cm}
        \begin{figure}
            \centering
            \includegraphics[scale=0.32]{nogrowth_reactors_3-7.pdf}
            \caption{Number of advanced reactors deployed.}
            \label{fig:reactors_3-7}
    \end{figure}
    \end{columns}
\end{frame}

\begin{frame}
    \frametitle{Enriched Uranium Mass}
    \begin{columns}
        \column[t]{5cm}
            \begin{itemize}
                \item Scenario 5 (MMR + VOYGR) requires the largest average mass of 
                      enriched uranium
                \item Scenario 2 (MMR) requires the largest mass of \gls{HALEU}
                \item Scenario 3 (Xe-100) requires the smallest mass of enriched 
                      uranium
                \item Scenario 5 (MMR + VOYGR) requires the smallest mass of \gls{HALEU}
                \item Result of fuel utilization of each of the advanced 
                      reactors
            \end{itemize}
        \column[t]{5cm}
        \vspace{-0.8cm}
            \begin{figure}
                \centering
                \includegraphics[scale=0.4]{nogrowth_AR_uranium.pdf}
                \caption{Annual avergae mass of enriched uranium required to fuel
                advanced reactors in Scenarios 2-7.}
                \label{fig:uranium}

        \end{figure}
    \end{columns}
    

\end{frame}

\begin{frame}
    \frametitle{Feed Uranium Requirements}
    \begin{columns}
        \column[t]{5cm}
            \begin{itemize}
                \item Scenario 2 (MMR) requires the largest average mass of feed uranium
                \item Scenario 5 (MMR + VOYGR) requires the next largest average
                \item The other scenarios require a feed uranium on a similar order 
                      of magnitude to \glspl{LWR}
                \item Feed uranium is a function of the mass of the product 
                      and the product assay
                
            \end{itemize}
        \column[t]{5cm}
        \vspace{-1cm}
        \begin{figure}
                \centering
                \includegraphics[height=0.4\textheight]{nogrowth_AR_feed.pdf}
                \caption{Annual average feed uranium required to produce enriched uranium 
                for advanced reactors in Scenarios 2-7.}
                \label{fig:feed}
        \end{figure}
    \end{columns}
\end{frame}

\begin{frame}
    \frametitle{\gls{SWU} Requirements}
    \begin{columns}
        \column[t]{5cm}
            \begin{itemize}
                \item Follows similar pattern to feed uranium  mass 
                \item Scenario 2 (MMR) requires the largest average \gls{SWU} 
                \item The other scenarios are comparable for the average 
                      capacity they require
                \item \gls{SWU} capacity is a function of product mass and 
                      product assay
                
            \end{itemize}
        \column[t]{5cm}
        \vspace{-1cm}
        \begin{figure}
                \centering
                \includegraphics[height=0.4\textheight]{nogrowth_AR_swu.pdf}
                \caption{Annual average \gls{SWU} capacity required to produce 
                enriched uranium for and advanced reactors in Scenarios 2-7.}
                \label{fig:swu}
        \end{figure}
    \end{columns}
\end{frame}

\begin{frame}
    \frametitle{\gls{SNF} Discharged}
    \begin{columns}
        \column[t]{5cm}
            \begin{itemize}
                \item Scenario 5 (MMR + VOYGR) discharges the largest mass of \gls{SNF}
                \item Scenario 3 (Xe-100) discharges the least \gls{SNF}
                \item Cumulative masses range between 25,654 MT (Scenario 3) to 
                      112,913 MT (Scenario 5)
                
            \end{itemize}
        \column[t]{5cm}
        \vspace{-1cm}
        \begin{figure}
                \centering
                \includegraphics[height=0.4\textheight]{nogrowth_AR_waste.pdf}
                \caption{Annual average mass of \gls{SNF} discharged from 
                advanced reactors in Scenarios 2-7.}
                \label{fig:waste}
        \end{figure}
    \end{columns}
\end{frame}